\chapter{Podsumowanie}


Zaprezentowana w niniejszym rozdziale warstwa użytkowa systemu „Wito’s” dowodzi, że zaprojektowany interfejs graficzny w pełni i z sukcesem realizuje wszystkie początkowe założenia funkcjonalne oraz biznesowe projektu. Architektura interfejsu została logicznie podzielona na dwie zintegrowane ze sobą strefy: front-end dedykowany klientom restauracji oraz kompleksowy back-office przeznaczony dla personelu zarządzającego.
\vspace{5pt}

Z perspektywy klienta końcowego, aplikacja dostarcza nowoczesnych i intuicyjnych narzędzi, które maksymalnie upraszczają proces zakupowy. Przejrzyste menu, błyskawiczny proces finalizacji zamówienia w koszyku oraz rozbudowany, interaktywny „Kreator pizzy” znacząco podnoszą jakość doświadczeń użytkownika (UX). Dodatkowo, osobisty portal „Moje konto” buduje zaufanie do marki, oddając w ręce klienta pełną kontrolę nad jego danymi, historią transakcji oraz procesem ewentualnych reklamacji.
\vspace{5pt}

Z punktu widzenia obsługi lokalu, wdrożone moduły administracyjne – obejmujące główny pulpit zarządzania asortymentem, cyfrowe stanowisko kuchni oraz mobilny panel dostawcy – całkowicie eliminują potrzebę stosowania tradycyjnych, podatnych na błędy papierowych rozwiązań. Pełna cyfryzacja obiegu informacji gwarantuje natychmiastową synchronizację danych w czasie rzeczywistym, co bezpośrednio przekłada się na krótszy czas realizacji zamówień, optymalizację logistyki i mniejszą liczbę pomyłek.
\vspace{5pt}

Warto również podkreślić innowacyjne podejście metodyczne, zastosowane podczas samej fazy prototypowania interfejsu. W celu zoptymalizowania czasu pracy nad warstwą wizualną, wykorzystano zaawansowane narzędzia ekosystemu Figma, w tym funkcje oparte na sztucznej inteligencji, takie jak narzędzie Figma Make (Make Designs). Zastosowanie tego rozwiązania pozwoliło na drastyczne przyspieszenie procesu generowania powtarzalnych widoków, ustandaryzowanych komponentów interfejsu oraz makiet strukturalnych (wireframes). Automatyzacja rutynowych zadań projektowych pozwoliła przenieść ciężar prac analitycznych na dopracowanie skomplikowanej logiki biznesowej, takiej jak ścieżka personalizacji w kreatorze produktu czy ergonomia paneli administracyjnych.
\vspace{5pt}

Kluczowym osiągnięciem zrealizowanego interfejsu jest również rygorystyczna implementacja standardów Responsive Web Design (RWD). Konsekwentna adaptacja każdego z opisywanych widoków do specyfiki ekranów dotykowych udowadnia, że system „Wito’s” jest rozwiązaniem wysoce elastycznym i gotowym do bezproblemowego działania na każdym urządzeniu – począwszy od domowego komputera klienta, aż po smartfon kuriera znajdującego się w trasie.

